%% ============================================================
%% Teil I: Python kompakt – was du für ROS2 wirklich brauchst
%% Band 2 – Kapitel 1 bis 5
%% ============================================================

\part{Python kompakt}

% ─────────────────────────────────────────────────────────────
\chapter{Wir wollen Python einsetzen – welche Version, welcher Stil?}
% ─────────────────────────────────────────────────────────────

\begin{lernziele}
Du kennst die Python-Version, die ROS2~Jazzy voraussetzt,
verstehst den Unterschied zwischen Systeminstallation und
virtueller Umgebung in ROS2-Projekten und kannst
Type~Hints, f-Strings und den \texttt{bytes}/\texttt{str}-Unterschied
bei pyserial erklären.
\end{lernziele}

\section{Situation}

\texttt{sensor\_bridge.py} aus Band~1 funktioniert.
Jetzt soll er zum ROS2-Node werden.
In \texttt{rclpy} sieht der erste Import so aus:

\begin{lstlisting}
import rclpy
from rclpy.node import Node
\end{lstlisting}

Das klappt -- aber nur mit der richtigen Python-Version
und nur, wenn ROS2 korrekt geladen ist.
Außerdem: Sobald du \texttt{def lese\_arduino(self) -> None:} schreibst
und VS~Code kein Autovervollständigung liefert, fragst du dich:
Was fehlt hier?

\begin{aufgabeBox}
Stelle sicher, dass Python in der richtigen Version läuft,
verstehe den \texttt{bytes}/\texttt{str}-Unterschied bei pyserial
und füge Type~Hints in \texttt{sensor\_bridge.py} ein.
\end{aufgabeBox}

\section{Erster Lösungsversuch}

\begin{lstlisting}
python --version
# Python 2.7.18 -- falsch!
python3 --version
# Python 3.12.3 -- richtig
\end{lstlisting}

\begin{problemBox}
ROS2~Jazzy braucht Python~3.10 oder neuer.
\texttt{python} (ohne "`3"') zeigt auf Python~2.7 auf manchen
Systemen -- und damit auf eine Version, die ROS2 nicht versteht.
Außerdem: rclpy ist systemweit installiert,
nicht in einer virtuellen Umgebung.
pip-Pakete in einem \texttt{venv} sehen rclpy nicht.
\end{problemBox}

\section{Was wir dafür brauchen}

\begin{wissenBox}
\textbf{Drei Grundregeln für Python in ROS2-Projekten:}

\begin{enumerate}
\item \textbf{Immer \texttt{python3}}, nie \texttt{python}.
  ROS2 sourced seinen Pfad via \texttt{source /opt/ros/jazzy/setup.bash} --
  danach zeigt \texttt{python3} garantiert auf die richtige Version.

\item \textbf{Kein \texttt{venv} für ROS2-Nodes.}
  \texttt{rclpy} ist systemweit installiert
  (\texttt{/opt/ros/jazzy/lib/python3.12/site-packages/rclpy}).
  Ein \texttt{venv} sieht das nicht.
  Für Drittbibliotheken (\texttt{pyserial}, \texttt{pyyaml}, \texttt{numpy}):
  entweder \texttt{pip3 install --break-system-packages}
  oder im Docker-Container wie in Band~1.

\item \textbf{Type~Hints aktivieren.}
  Pylance in VS~Code liest Type~Hints und zeigt dir Fehler
  schon beim Schreiben -- bevor der Node läuft.
\end{enumerate}
\end{wissenBox}

\subsection{bytes vs. str – der pyserial-Unterschied}

Pyserial liefert immer \texttt{bytes}, nie \texttt{str}:

\begin{lstlisting}
import serial
ser = serial.Serial('/dev/arduino', 9600, timeout=1)

zeile = ser.readline()
# type: bytes
# Wert: b'R=87 cm  L=124 cm\r\n'

text = zeile.decode('utf-8').strip()
# type: str
# Wert: 'R=87 cm  L=124 cm'
\end{lstlisting}

Serielle Kommunikation ist immer binär -- der Port weiß nicht,
ob das ASCII, UTF-8 oder Rohdaten sind.
\texttt{.decode('utf-8')} ist deshalb nicht optional,
sondern Pflicht.

\subsection{Type Hints -- der freiwillige Vertrag}

\begin{lstlisting}
# Ohne Type Hints: funktioniert, aber unklar
def parse_zeile(zeile):
    pass

# Mit Type Hints: VS Code kennt Ein- und Ausgabe
def parse_zeile(zeile: str) -> tuple[int, int] | None:
    """Parst 'R=87 cm  L=124 cm' zu (87, 124) oder None."""
    pass
\end{lstlisting}

Type~Hints ändern das Laufzeitverhalten nicht.
Aber Pylance liest sie -- und zeigt dir sofort, wenn du
\texttt{parse\_zeile(87)} statt \texttt{parse\_zeile("R=87 cm")} schreibst.

\section{Schritt-für-Schritt-Lösung}

\begin{loesungBox}
\textbf{Schritt~1:} Python-Version und ROS2 prüfen.
\begin{lstlisting}[language=bash]
source /opt/ros/jazzy/setup.bash
python3 --version  # muss 3.10+ sein
python3 -c "import rclpy; print('rclpy OK')"
\end{lstlisting}

\textbf{Schritt~2:} \texttt{sensor\_bridge.py} mit Type~Hints versehen.
\begin{lstlisting}
#!/usr/bin/env python3
# stufe2_ros2/src/sensor_bridge.py -- Version 0.4
import re
import serial
import yaml
from pathlib import Path

def parse_zeile(zeile: str) -> tuple[int, int] | None:
    """Parst 'R=87 cm  L=124 cm' -> (87, 124) oder None."""
    m = re.match(r'R=(\d+)\s+cm\s+L=(\d+)\s+cm', zeile.strip())
    if m:
        return int(m.group(1)), int(m.group(2))
    return None

def cm_zu_meter(cm: int) -> float:
    """Einheitenumrechnung für sensor_msgs/Range."""
    return float(cm) / 100.0
\end{lstlisting}

\textbf{Schritt~3:} Type-Check mit mypy.
\begin{lstlisting}[language=bash]
pip3 install mypy --break-system-packages
mypy src/sensor_bridge.py
# Success: no issues found in 1 source file
\end{lstlisting}
\end{loesungBox}

\begin{projektBox}[robotik-uno-q -- Projektstand]
\textbf{Bisher:} \texttt{sensor\_bridge.py} aus Band~1 -- ohne Type Hints.\\
\textbf{Heute:} Type Hints für alle Funktionen, \texttt{cm\_zu\_meter()} hinzugefügt,
mypy bestätigt Korrektheit.\\
\textbf{Nächstes Kapitel:} Den Bridge als Klasse aufbauen -- Vorbereitung auf rclpy.
\end{projektBox}

\begin{merksatzBox}
\textbf{Merksatz:} Immer \texttt{python3}. Kein \texttt{venv} für ROS2-Nodes.
\texttt{bytes.decode('utf-8')} ist Pflicht bei pyserial.
Type~Hints kosten nichts und geben dir Pylance als kostenlosen Code-Reviewer.
\end{merksatzBox}

\begin{fehlerBox}
\begin{itemize}
\item \texttt{python} statt \texttt{python3}: Führt zu Python~2, rclpy-Import schlägt fehl.
\item \texttt{venv} mit rclpy: \texttt{ModuleNotFoundError: No module named 'rclpy'} --
  rclpy ist nicht im venv sichtbar.
\item \texttt{int("87.5")}: \texttt{ValueError} -- erst \texttt{float()}, dann \texttt{int()}.
\item \texttt{tuple[int, int]} ohne Python~3.9+: \texttt{from typing import Tuple} verwenden.
\end{itemize}
\end{fehlerBox}

\begin{uebungBox}
\begin{enumerate}
\item Füge Type~Hints in die gesamte \texttt{sensor\_bridge.py} aus Band~1 ein
  und führe \texttt{mypy} aus. Behebe alle Warnungen.
\item Schreibe eine Funktion \texttt{validiere\_abstand(cm: int) -> bool},
  die \texttt{True} zurückgibt, wenn der Wert zwischen 2 und 400\,cm liegt.
\item Was passiert, wenn der Arduino \texttt{"Bremsen\textbackslash r\textbackslash n"} sendet
  statt \texttt{"R=87 cm  L=124 cm"}? Passe \texttt{parse\_zeile()} entsprechend an.
\end{enumerate}
\end{uebungBox}

\begin{quizBox}
\begin{enumerate}
\item Warum liefert pyserial \texttt{bytes} statt \texttt{str}?
\item Was ändert ein Type~Hint am Laufzeitverhalten?
\item Warum ist ein \texttt{venv} für ROS2-Nodes problematisch?
\end{enumerate}
\textit{Antworten:}
1.~Serielle Kommunikation ist binär; der Typ des Inhalts ist dem Port unbekannt.
2.~Nichts -- Type~Hints sind reine Metadaten für Werkzeuge wie mypy und Pylance.
3.~rclpy ist systemweit unter \texttt{/opt/ros/} installiert und im venv nicht sichtbar.
\end{quizBox}

\begin{haubeBox}
\textbf{Warum PEP~484 -- Type~Hints erst 2014?}

Python wurde als dynamische Skriptsprache entworfen -- schnell zu schreiben,
kein Compiler nötig. Das war 1991 ein Vorteil.
Als Python-Projekte aber größer wurden (Django, NumPy, ROS2),
wurde das fehlende Typsystem zum Problem:
Fehler zeigten sich erst zur Laufzeit.

PEP~484 löst das ohne den Interpreter zu ändern:
Type~Hints sind optionale Metadaten, die statische Analysetools lesen.
Der Interpreter ignoriert sie vollständig.
Das ist ein typischer Python-Kompromiss -- Rückwärtskompatibilität
über alles.
\end{haubeBox}

\begin{robotikBox}
In \texttt{sensor\_msgs/Range} ist das \texttt{range}-Feld ein \texttt{float32}.
\texttt{cm\_zu\_meter()} liefert genau das.
In Teil~III siehst du, wie der Publisher-Node das Range-Objekt befüllt --
und warum \texttt{float32} statt \texttt{float64} die richtige Wahl für Sensordaten ist.
\end{robotikBox}


% ─────────────────────────────────────────────────────────────
\chapter{Wir wollen einen ROS2-Node als Klasse schreiben –
  was bedeuten \texttt{self} und \texttt{super()}?}
% ─────────────────────────────────────────────────────────────

\begin{lernziele}
Du verstehst Klassen, Instanzmethoden und Vererbung in Python,
kannst erklären, warum jeder rclpy-Node eine Klasse ist,
und schreibst deinen ersten \texttt{MinimalNode},
der sich korrekt initialisiert.
\end{lernziele}

\section{Situation}

Die rclpy-Dokumentation zeigt diesen Code:

\begin{lstlisting}
class MinimalNode(Node):
    def __init__(self):
        super().__init__('minimal_node')
\end{lstlisting}

Drei Dinge, die Anfänger verwirren:
Was ist \texttt{Node} -- und warum erbt MinimalNode davon?
Was bedeutet \texttt{self}?
Warum \texttt{super().\_\_init\_\_()}?

\begin{aufgabeBox}
Verstehe das Klassensystem von Python so weit,
dass du jeden rclpy-Node selbst schreiben kannst --
ohne die Struktur auswendig zu lernen.
\end{aufgabeBox}

\section{Das Klassensystem in Python}

\begin{wissenBox}
\textbf{Eine Klasse ist eine Vorlage. Eine Instanz ist das konkrete Objekt.}

\begin{lstlisting}
class Sensor:
    """Vorlage für einen Ultraschallsensor."""

    def __init__(self, name: str, max_abstand: float) -> None:
        # self = dieses konkrete Objekt (die Instanz)
        self.name = name                # Instanzvariable
        self.max_abstand = max_abstand  # Instanzvariable
        self._messwerte: list[float] = []  # privat, per Konvention

    def messe(self, rohwert: int) -> float:
        """Konvertiert Rohwert und speichert ihn."""
        wert = float(rohwert) / 100.0
        self._messwerte.append(wert)
        return wert

    def mittelwert(self) -> float:
        """Gibt den Mittelwert aller Messungen zurück."""
        if not self._messwerte:
            return 0.0
        return sum(self._messwerte) / len(self._messwerte)

# Zwei unabhängige Instanzen
sensor_rechts = Sensor("rechts", max_abstand=4.0)
sensor_links  = Sensor("links",  max_abstand=4.0)

sensor_rechts.messe(87)   # -> 0.87
sensor_links.messe(124)   # -> 1.24
\end{lstlisting}

\textbf{Merke:} \texttt{self} ist kein Schlüsselwort -- es ist
eine Konvention. Python übergibt die Instanz automatisch
als erstes Argument jeder Methode.
\texttt{sensor\_rechts.messe(87)} ruft intern
\texttt{Sensor.messe(sensor\_rechts, 87)} auf.
\end{wissenBox}

\subsection{Vererbung -- \texttt{Node} als Basisklasse}

\texttt{rclpy.node.Node} stellt alles bereit, was ein ROS2-Node braucht:
Logger, Publisher, Subscriber, Timer, Parameter.
Dein Node erbt das alles, ohne es neu zu schreiben.

\begin{lstlisting}
from rclpy.node import Node

class ArduinoBridgeNode(Node):
    def __init__(self) -> None:
        # 1. Basisklasse initialisieren -- MUSS als Erstes kommen
        super().__init__('arduino_bridge')
        # Jetzt sind self.get_logger(), self.create_publisher() usw. verfügbar

        # 2. Eigene Initialisierung
        self._port: str = self.declare_parameter(
            'port', '/dev/arduino'
        ).get_parameter_value().string_value

        self.get_logger().info(f"ArduinoBridgeNode gestartet. Port: {self._port}")
\end{lstlisting}

\texttt{super().\_\_init\_\_('arduino\_bridge')} tut zwei Dinge:
\begin{enumerate}
\item Es initialisiert den ROS2-Node mit dem Namen \texttt{'arduino\_bridge'}.
  Dieser Name erscheint in \texttt{ros2 node list}.
\item Es richtet alle internen ROS2-Strukturen ein, die Node braucht.
  Ohne diesen Aufruf stürzt jeder Zugriff auf \texttt{self.get\_logger()} ab.
\end{enumerate}

\section{Schritt-für-Schritt-Lösung}

\begin{loesungBox}
Erstelle \texttt{stufe2\_ros2/src/arduino\_bridge\_node.py} als Klasse:

\begin{lstlisting}
#!/usr/bin/env python3
# stufe2_ros2/src/arduino_bridge_node.py -- Version 0.1
"""ArduinoBridgeNode: liest Serial, published Sensordaten."""

import rclpy
from rclpy.node import Node

class ArduinoBridgeNode(Node):
    """Liest HC-SR04-Daten vom Arduino und published sie als ROS2-Topics."""

    def __init__(self) -> None:
        super().__init__('arduino_bridge')
        self.get_logger().info("ArduinoBridgeNode bereit.")

def main(args: list[str] | None = None) -> None:
    rclpy.init(args=args)
    node = ArduinoBridgeNode()
    try:
        rclpy.spin(node)
    except KeyboardInterrupt:
        pass
    finally:
        node.destroy_node()
        rclpy.shutdown()

if __name__ == '__main__':
    main()
\end{lstlisting}

Testen (ohne ROS2-Paket, direkt):
\begin{lstlisting}[language=bash]
source /opt/ros/jazzy/setup.bash
python3 src/arduino_bridge_node.py
# [INFO] [arduino_bridge]: ArduinoBridgeNode bereit.
# Strg+C zum Beenden
\end{lstlisting}
\end{loesungBox}

\begin{projektBox}[robotik-uno-q -- Projektstand]
\textbf{Bisher:} \texttt{sensor\_bridge.py} als Funktion.\\
\textbf{Heute:} \texttt{arduino\_bridge\_node.py} als Klasse -- bereit für Publisher.\\
\textbf{Nächstes Kapitel:} Type~Hints und Datenstrukturen vertiefen.
\end{projektBox}

\begin{merksatzBox}
\textbf{Merksatz:} Jeder rclpy-Node ist eine Klasse, die von \texttt{Node} erbt.
\texttt{super().\_\_init\_\_('node\_name')} muss als Erstes in \texttt{\_\_init\_\_()} stehen.
Der Node-Name erscheint in \texttt{ros2 node list} -- er sollte eindeutig
und beschreibend sein.
\end{merksatzBox}

\begin{fehlerBox}
\begin{itemize}
\item \texttt{super().\_\_init\_\_()} vergessen:
  \texttt{AttributeError: 'ArduinoBridgeNode' has no attribute '\_node'} --
  der Node ist nicht initialisiert.
\item Node-Name mit Slash: \texttt{'/arduino\_bridge'} ist falsch;
  der Leading-Slash kommt vom Namespace, nicht vom Namen selbst.
\item \texttt{rclpy.init()} vergessen: \texttt{RuntimeError: rclpy not initialized}.
\end{itemize}
\end{fehlerBox}

\begin{haubeBox}
\textbf{Warum ist \texttt{super()} in Python so merkwürdig?}

In C++ schreibt man \texttt{Node::Node(name)} explizit mit dem Klassennamen.
Python's \texttt{super()} ist dynamisch: Es findet die richtige Basisklasse
anhand der sogenannten MRO (Method Resolution Order) -- dem Algorithmus,
der bei Mehrfachvererbung bestimmt, welche Methode zuerst aufgerufen wird.
Für einfache Vererbung wie \texttt{ArduinoBridgeNode(Node)} ist \texttt{super()}
identisch mit \texttt{Node.\_\_init\_\_(self, 'arduino\_bridge')}.
Die MRO ist nur bei \texttt{class C(A, B):} relevant.
\end{haubeBox}


% ─────────────────────────────────────────────────────────────
\chapter{Wir wollen Fehler abfangen – \texttt{try/except} und Context Manager}
% ─────────────────────────────────────────────────────────────

\begin{lernziele}
Du kannst \texttt{SerialException} und andere Hardware-Fehler sauber abfangen,
verstehst Context~Manager (\texttt{with}-Statement) und
verwendest den Python-Logger statt \texttt{print()}.
\end{lernziele}

\section{Situation}

\texttt{arduino\_bridge\_node.py} stürzt ab,
sobald der Arduino nicht angesteckt ist:

\begin{lstlisting}
ser = serial.Serial('/dev/arduino', 9600)
# serial.serialutil.SerialException: [Errno 2] No such file or directory: '/dev/arduino'
\end{lstlisting}

Ein ROS2-Node darf nicht abstürzen -- er soll warten, melden und neu versuchen.

\begin{aufgabeBox}
Mache \texttt{arduino\_bridge\_node.py} robust gegen fehlende Hardware:
SerialException abfangen, sauber loggen, Reconnect-Logik hinzufügen.
\end{aufgabeBox}

\section{try/except/finally in Python}

\begin{wissenBox}
\begin{lstlisting}
try:
    # Code, der einen Fehler auslösen könnte
    ser = serial.Serial('/dev/arduino', 9600, timeout=1)
except serial.SerialException as e:
    # Fehler abfangen -- e enthält die Fehlermeldung
    self.get_logger().error(f"Serial-Fehler: {e}")
except OSError as e:
    # Breiter: alle OS-Fehler (Permission denied, etc.)
    self.get_logger().error(f"OS-Fehler: {e}")
else:
    # Wird nur ausgeführt, wenn KEIN Fehler aufgetreten ist
    self.get_logger().info("Verbunden.")
finally:
    # Wird IMMER ausgeführt -- ideal für Aufräumarbeiten
    pass
\end{lstlisting}

\textbf{Faustregel:} Fange immer den spezifischsten Fehler zuerst.
\texttt{except Exception:} ist der letzte Ausweg -- es versteckt Bugs.
\end{wissenBox}

\subsection{Context Manager -- \texttt{with} statt \texttt{try/finally}}

\begin{lstlisting}
# Ohne Context Manager: Fehleranfällig
ser = serial.Serial('/dev/arduino', 9600)
try:
    data = ser.readline()
finally:
    ser.close()  # Wird vergessen? Port bleibt offen.

# Mit Context Manager: Sauber und sicher
with serial.Serial('/dev/arduino', 9600) as ser:
    data = ser.readline()
# ser.close() wird automatisch aufgerufen -- auch bei Ausnahmen
\end{lstlisting}

\begin{lstlisting}
# Robuste Reconnect-Logik für den Node
def _verbinde(self) -> bool:
    """Versucht, den Serial-Port zu öffnen. Gibt True zurück bei Erfolg."""
    try:
        self._ser = serial.Serial(self._port, self._baudrate, timeout=1.0)
        self.get_logger().info(f"Verbunden: {self._port}")
        return True
    except serial.SerialException as e:
        self.get_logger().warning(f"Kein Arduino: {e}. Neuer Versuch in 5s.")
        return False
\end{lstlisting}

\begin{projektBox}[robotik-uno-q -- Projektstand]
\textbf{Heute:} Robuste Fehlerbehandlung im \texttt{arduino\_bridge\_node}. \\
\textbf{Nächstes Kapitel:} Callbacks und Lambda -- wie ROS2-Timer funktionieren.
\end{projektBox}

\begin{merksatzBox}
\textbf{Merksatz:} \texttt{with serial.Serial(...) as ser:} ist besser als
manuellem \texttt{try/finally}. Spezifische Exceptions (\texttt{SerialException})
vor allgemeinen (\texttt{Exception}) fangen. ROS2-Nodes loggen mit
\texttt{self.get\_logger().info/warning/error()} -- nie mit \texttt{print()}.
\end{merksatzBox}

\begin{fehlerBox}
\begin{itemize}
\item \texttt{except Exception:} ohne spezifischen Typ: Versteckt echte Bugs
  (z.\,B. \texttt{TypeError} in der Verarbeitungslogik).
\item Port nicht schließen: Beim Neustart des Nodes
  ist der Port noch belegt (\texttt{Permission denied}).
\item \texttt{print()} statt Logger: Nachrichten erscheinen nicht in
  \texttt{ros2 log} und nicht in \texttt{journalctl}.
\end{itemize}
\end{fehlerBox}


% ─────────────────────────────────────────────────────────────
\chapter{Wir wollen Daten strukturieren –
  Dicts, Dataclasses und NamedTuples}
% ─────────────────────────────────────────────────────────────

\begin{lernziele}
Du kennst den Unterschied zwischen \texttt{dict},
\texttt{dataclass} und \texttt{NamedTuple} und wählst
das richtige Werkzeug für Sensordaten im Node.
\end{lernziele}

\section{Situation}

\texttt{parse\_zeile()} gibt \texttt{tuple[int, int]} zurück.
Das reicht für zwei Werte -- aber was passiert,
wenn der Arduino bald auch Batterieladung und Temperatur sendet?
\texttt{tuple[int, int, float, int]} wird unleserlich.

\begin{aufgabeBox}
Erstelle eine Datenstruktur \texttt{SensorDaten}, die Abstandswerte,
Zeitstempel und Statusinformationen klar bündelt.
\end{aufgabeBox}

\section{Die drei Optionen}

\begin{wissenBox}
\textbf{dict:} Flexibel, aber unsicher.
\begin{lstlisting}
data = {"rechts": 87, "links": 124, "einheit": "cm"}
data["rechts"]   # OK
data["Rechts"]   # KeyError -- Tippfehler zur Laufzeit
\end{lstlisting}

\textbf{NamedTuple:} Unveränderlich, typisiert, leichtgewichtig.
\begin{lstlisting}
from typing import NamedTuple
class SensorDaten(NamedTuple):
    rechts_cm: int
    links_cm:  int
    zeitstempel: float

daten = SensorDaten(rechts_cm=87, links_cm=124, zeitstempel=1720000000.0)
daten.rechts_cm   # 87 -- typsicher, kein KeyError
daten.rechts_cm = 90  # AttributeError -- unveränderlich
\end{lstlisting}

\textbf{dataclass:} Veränderlich, typisiert, mehr Kontrolle.
\begin{lstlisting}
from dataclasses import dataclass, field
import time

@dataclass
class SensorDaten:
    rechts_cm: int
    links_cm:  int
    zeitstempel: float = field(default_factory=time.time)
    gueltig: bool = True

daten = SensorDaten(rechts_cm=87, links_cm=124)
daten.rechts_cm = 90   # OK -- veränderlich
\end{lstlisting}
\end{wissenBox}

\textbf{Empfehlung für robotik-uno-q:}
\texttt{dataclass} für \texttt{SensorDaten},
weil Zeitstempel automatisch gesetzt werden soll
und Gültigkeitsflag nachträglich gesetzt werden kann.

\begin{lstlisting}
# stufe2_ros2/src/sensor_daten.py
from dataclasses import dataclass, field
import time

@dataclass
class SensorDaten:
    """Bündelt einen Messwert vom HC-SR04-Dual-Sensor."""
    rechts_cm:   int
    links_cm:    int
    zeitstempel: float = field(default_factory=time.time)
    gueltig:     bool  = True

    @property
    def minimum_cm(self) -> int:
        """Gibt den kleineren der beiden Abstandswerte zurück."""
        return min(self.rechts_cm, self.links_cm)

    @property
    def rechts_m(self) -> float:
        return self.rechts_cm / 100.0

    @property
    def links_m(self) -> float:
        return self.links_cm / 100.0
\end{lstlisting}

\begin{projektBox}[robotik-uno-q -- Projektstand]
\textbf{Heute:} \texttt{SensorDaten}-Dataclass erstellt.
Parser gibt jetzt \texttt{SensorDaten} zurück statt \texttt{tuple[int, int]}.\\
\textbf{Nächstes Kapitel:} Callbacks und Lambda für Timer.
\end{projektBox}

\begin{merksatzBox}
\textbf{Merksatz:}
\texttt{dict} für unbekannte Schlüssel zur Laufzeit (YAML-Config).
\texttt{NamedTuple} für unveränderliche Datenpakete.
\texttt{dataclass} für Objekte mit Methoden und Defaultwerten.
In rclpy-Nodes ist \texttt{dataclass} fast immer die richtige Wahl.
\end{merksatzBox}


% ─────────────────────────────────────────────────────────────
\chapter{Wir wollen Funktionen als Parameter übergeben –
  Callbacks und Lambda}
% ─────────────────────────────────────────────────────────────

\begin{lernziele}
Du verstehst, warum rclpy-Timer und Subscriber Funktionen
als Parameter erwarten, und kannst Callbacks,
Lambda-Ausdrücke und \texttt{functools.partial} korrekt einsetzen.
\end{lernziele}

\section{Situation}

\texttt{create\_timer(0.1, callback)} erwartet eine Funktion als zweites Argument.
Was genau übergibst du da?
Und was, wenn der Callback einen Parameter braucht,
den \texttt{create\_timer()} nicht kennt?

\begin{aufgabeBox}
Schreibe einen Timer-Callback für \texttt{arduino\_bridge\_node},
der alle 100\,ms die Serial-Leitung abfragt --
und verstehe, warum Funktionen in Python Objekte sind.
\end{aufgabeBox}

\section{Funktionen als Objekte}

\begin{wissenBox}
In Python ist eine Funktion ein Objekt -- sie kann in einer Variable
gespeichert, als Parameter übergeben und zur Laufzeit erstellt werden.

\begin{lstlisting}
def lese_sensor() -> None:
    print("Sensor wird gelesen...")

# Funktion als Objekt -- ohne ()
aktion = lese_sensor        # aktion IST die Funktion
aktion()                    # jetzt wird sie aufgerufen

# Lambda: anonyme Einzeiler-Funktion
quadrat = lambda x: x * x
quadrat(5)                  # -> 25

# functools.partial: Funktion mit vorausgefüllten Argumenten
import functools
def logge(level: str, nachricht: str) -> None:
    print(f"[{level}] {nachricht}")

logge_info = functools.partial(logge, "INFO")
logge_info("Sensor verbunden.")   # -> [INFO] Sensor verbunden.
\end{lstlisting}
\end{wissenBox}

\subsection{Timer-Callback in rclpy}

\begin{lstlisting}
class ArduinoBridgeNode(Node):
    def __init__(self) -> None:
        super().__init__('arduino_bridge')

        # create_timer(periode_in_sek, callback_funktion)
        # Alle 100 ms wird self._lese_arduino aufgerufen
        self._timer = self.create_timer(0.1, self._lese_arduino)

    def _lese_arduino(self) -> None:
        """Wird alle 100 ms vom Timer aufgerufen."""
        # Hier: Serial-Zeile lesen, parsen, publishen
        self.get_logger().debug("Timer-Tick")
\end{lstlisting}

\texttt{self.\_lese\_arduino} ist ein \emph{Bound Method} --
eine Methode, die bereits an die Instanz gebunden ist.
Wenn der Timer sie aufruft, hat \texttt{self} seinen Wert.

\begin{lstlisting}
# Lambda für schnelle Inline-Callbacks (bei einfachen Fällen)
self._timer = self.create_timer(
    1.0,
    lambda: self.get_logger().info("Heartbeat")
)
\end{lstlisting}

\begin{projektBox}[robotik-uno-q -- Projektstand]
\textbf{Heute:} Timer-Callback für 10\,Hz-Sensor-Polling implementiert.\\
\textbf{Nächster Teil:} Robotik-Bibliotheken -- NumPy, Matplotlib, OpenCV.
\end{projektBox}

\begin{merksatzBox}
\textbf{Merksatz:} Funktionen ohne \texttt{()} sind Objekte -- übergib sie so.
\texttt{self.\_methode} ist bereits an \texttt{self} gebunden.
Lambda für Einzeiler, \texttt{functools.partial} für Funktionen
mit vorausgefüllten Argumenten.
\end{merksatzBox}

\begin{fehlerBox}
\begin{itemize}
\item \texttt{create\_timer(0.1, self.\_lese\_arduino())}:
  Die Klammern rufen die Methode sofort auf und übergeben den
  Rückgabewert (\texttt{None}) als Callback -- kein Timer-Tick.
\item Lambda mit Seiteneffekten in Schleifen:
  \texttt{for i in range(3): callbacks.append(lambda: print(i))} --
  alle drei Lambdas drucken den letzten Wert von \texttt{i}.
  Lösung: \texttt{lambda i=i: print(i)}.
\end{itemize}
\end{fehlerBox}
